\documentclass[12pt]{ociamthesis}  % default square logo 
%\documentclass[12pt,beltcrest]{ociamthesis} % use old belt crest logo
%\documentclass[12pt,shieldcrest]{ociamthesis} % use older shield crest logo

%load any additional packages
\usepackage{amssymb}
\usepackage{caption}
\usepackage{multirow}
\usepackage{listings}
\usepackage{longtable}

%input macros (i.e. write your own macros file called mymacros.tex 
%and uncomment the next line)
%\include{mymacros}

\title{\large INTEGRASI DATA INBOUND PADA SAAT PENGUJIAN SMART CONVEYOR DAN SISTEM WAREHOUSE MANAGEMENT SYSTEM  \\[5ex]   %your thesis title,
\small{Laporan ini dibuat untuk memenuhi persyaratan kelulusan matakuliah\\[-4ex]
	Proyek 3}} %note \\[1ex] is a line break in the title

\author{\large\textbf{Alit Fajar Kurniawan}}             %your name
\college{\large\textbf{1.17.4.057}\\[5ex]}

%\textit{In Partial Fulfilment of The Requirements for The Degree of Applied Bachelor of Informatics Engineering}}  %your college

%\renewcommand{\submittedtext}{change the default text here if needed}
\degree{\textbf{POLITEKNIK POS INDONESIA}}     %the degree
%\degreedate{Bandung}         %the degree date
\degreedate{\textbf{2019}}  

%end the preamble and start the document
\begin{document}

%this baselineskip gives sufficient line spacing for an examiner to easily
%markup the thesis with comments
\baselineskip=18pt plus1pt

%set the number of sectioning levels that get number and appear in the contents
\setcounter{secnumdepth}{12}
\setcounter{tocdepth}{12}

\maketitle                  % create a title page from the preamble info
\begin{romanpages}
\addcontentsline{toc}{chapter}{COVER}
\begin{dedication}

`Jika Kamu tidak dapat menahan lelahnya belajar, \\
Maka kamu harus sanggup menahan perihnya Kebodohan.'\\ 
~Imam Syafi'i~\\

\end{dedication}   % include a dedication.tex file

\addcontentsline{toc}{chapter}{LEMBAR PENGESAHAN LAPORAN}
\begin{pengesahan}
\thispagestyle{plain}
	
\begin{figure}[!htbp]
	\centering
	\includegraphics[width=32mm]{figures/logo-poltekpos} \\[4ex]
\end{figure} 

\centering\textbf{LEMBAR PENGESAHAN LAPORAN PROYEK 3} \\[2ex]

\centering{\textbf{Disusun Oleh,}} \\
\centering{\textbf{Alit Fajar Kurniawan}} \\
\centering{\textbf{1.17.4.057}} \\[2ex]

Bandung, 27 Agustus 2019 \\[2ex]

\begin{table}[h]
	\begin{tabular}{ccc}
		Dosen Pembimbing Utama,                  &                      & Dosen Pembimbing Pendamping,       \\
		&                      &                                    \\
		&                      &                                    \\
		&                      &                                    \\
		&                      &                                    \\
		\underline{Rolly Maulana Awangga, S.T., M.T.}        &                      & \underline{Syafrial Fachri Pane, S.T., M.T.I.} \\
		NIK: 117.86.219                          &                      & NIK: 117.88.233                    \\
		\multicolumn{1}{l}{}                     & \multicolumn{1}{l}{} & \multicolumn{1}{l}{}               \\
		Mengetahui,                              &                      &                                    \\
		Kaprodi D4 Teknik Informatika            &                      & Dosen Penguji,                     \\
		&                      &                                    \\
		&                      &                                    \\
		&                      &                                    \\
		&                      &                                    \\
		\underline{M. Yusril Helmi Setyawan, S.Kom., M.Kom.} &                      & \underline{Rd. Nuraini, Siti Fatonah, S.S., M.Hum.}    \\
		NIK: 113.74.163                          &                      & NIK: 217.72.187                   
	\end{tabular}
\end{table}

\end{pengesahan}     %include a pengesahan.tex file

\addcontentsline{toc}{chapter}{LEMBAR PERNYATAAN}
\sloppy

\begin{pernyataan}
\thispagestyle{plain}

Dengan ini saya,

\begin{table}[ht]
	\centering
	\begin{tabular}{lll}
		Nama     & : & Alit Fajar Kurniawan                                                                                                                             \\
		NPM       & : & 1.17.4.057                                                                                                                                    \\
		Judul TA & : & \begin{tabular}[c]{@{}l@{}}INTEGRASI DATA INBOUND PADA SAAT PENGUJIAN  \\ SMART CONVEYOR DAN SISTEM WAREHOUSE MANAGEMENT SYSTEM \end{tabular}
	\end{tabular}
\end{table}

Memberikan kepada Politeknik Pos Indonesia hak non-eksklusif untuk menyimpan, memperbanyak, dan menyebarluaskan tesis karya saya, secara keseluruhan atau hanya sebagian atau hanya ringkasan saja, dalam bentuk format tercetak dana tau elektronik. \\[2ex]
Menyatakan bahwa saya, akan mempertahankan hak exclusive saya, untuk menggunakan seluruh atau sebagian isi tesis saya, guna pengembangan karya di masa depan, misalnya bentuk artikelm buku, perangkat lunak, ataupun sistem informasi. \\[2ex]
\textit{Hereby grant to my school, Politeknik Pos Indonesia the non-exclusive right to archive, reproduce, and distribute my thesis, in whole or in part whether in the form of printed and electronic formats.} \\[2ex]
\textit{I acknowledge that I retain exclusive rights of my thesis by using all or part of it in the future work or outputs, such as article, book software and information system.} \\[1ex]
\begin{flushright}
Bandung, 27 Agustus 2019 \\[10ex]


Maulyanda
\end{flushright}

\end{pernyataan}    %include a pernyataan.tex file

\addcontentsline{toc}{chapter}{HALAMAN PERNYATAAN}
\sloppy

\begin{halaman}
\thispagestyle{plain}

Saya, nama Mahasiwa, NIM 1154008 menyatakan dengan sebenar-benarnya bahwa Tugas Akhir/Skripsi saya berjudul ``JUDUL LAPORAN TUGAS AKHIR'' adalah merupakan gagasan dan hasil penelitian saya sendiri dengan bimbingan Dosen Pembimbing.
Saya juga menyatakan dengan sebenarnya bahwa isi Tugas Akhir/Skripsi ini tidak merupakan jiplakan dan bukan pula dari karya orang lain, kecuali kutipan dari literatur dana tau hasil wawancara tertulis yang saya acu dan telah saya sebutkan di Daftar Acuan dan Daftar Pustaka.
Demikian pernyataan ini saya buat dengan sebenarnya dan saya bersedia menerima sanksi apabila ternyata pernyataan saya ini tidak benar. \\[2ex]
\textit{I, Name Mahaiswa, Student ID 1154008 truly acknowledge that my thesis with title ``JUDUL LAPORAN TUGAS AKHIR'' is my concept and project result with guidance from supervisor.
I, also truly acknowledge that content of this thesis are not copyed and not from another people work, except my citation from literature or written interview result and already write in reference list and bibliography list. 
That’s my acknowledge were truly made and if in reality this acknowledge weren’t true, I willing sanction.}
\begin{flushright}
Bandung, 27 Agustus 2019\\
Yang menyatakan \\[11ex]

Nama Mahasiswa \\
NPM
\end{flushright}

\end{halaman}      %include a halaman.tex file

\addcontentsline{toc}{chapter}{ABSTRAK}
\begin{abstrak}
Buku Pedoman ini dibuat dengan tujuan memberikan acuan, bagi mahasiswa Tingkat Akhir dan dosen
Pembimbing. Pada intinya buku ini menjelaskan secara lengkap tentang Standar pengerjaan Intership  dan 
Tugas Akhir
di Program Studi D4 Teknik Informatika, dan juga mengatur mekanisme, teknik penulisan, serta
penilaiannya.Dengan demikian diharapkan semua pihak yang terlibat dalam aktivitas Bimbingan Mahasiswa Tingkat Akhir
berjalan lancar dan sesuai dengan standar.
\end{abstrak}  % include the abstract

\addcontentsline{toc}{chapter}{\textit{ABSTRACT}}
\begin{abstract}
Buku Pedoman ini dibuat dengan tujuan memberikan acuan, bagi mahasiswa Tingkat Akhir dan dosen
Pembimbing. Pada intinya buku ini menjelaskan secara lengkap tentang Standar pengerjaan Intership  dan 
Tugas Akhir
di Program Studi D4 Teknik Informatika, dan juga mengatur mekanisme, teknik penulisan, serta
penilaiannya.Dengan demikian diharapkan semua pihak yang terlibat dalam aktivitas Bimbingan Mahasiswa Tingkat Akhir
berjalan lancar dan sesuai dengan standar.
\end{abstract}

% start roman page numbering
\addcontentsline{toc}{chapter}{KATA PENGANTAR}
\begin{acknowledgements}
\thispagestyle{plain}
Assalamualaikum warahmatullahi wabarakatuh. Segala puji bagi Allah SWT yang telah memberikan kemudahan sehingga dapat menyelesaikan laporan Tugas Akhir ini, tanpa pertolongan-Nya mungkin penulis tidak akan sanggup menyelesaikannya dengan baik. Shalawat dan salam semoga terlimpah curahkan kepada Nabi Muhammad SAW beserta sahabat dan keluarga Beliau.

Laporan ini disusun untuk memenuhi kelulusan matakuliah Tugas Akhir pada Program Studi DIV Teknik Informatika. Proses Tugas Akhir ini juga tidak terlepas dari bantuan berbagai pihak. Oleh karena itu, pada kata pengantar ini penulis menyampaikan teriamakasih kepada :
\begin{enumerate}

\item Rolly Maulana Awangga, S.T., M.T. selaku Pembimbing Utama dalam penyusunan laporan Tugas Akhir ini;
\item	Syafrial Fachri Pane, S.T., M.T.I. selaku Pembimbing Pendamping dalam penyusunan laporan Tugas Akhir dan Koordinator Tugas Akhir Tahun Akademik 2018/2019;
\item	M. Yusril Helmi Setyawan, S.Kom., M.Kom. selaku Ketua Program Studi DIV Teknik Informatika Tahun Akademik 2018/2019;
\item	Dr. Ir. Agus Purnomo, M.T. selaku Direktur Politeknik Pos Indonesia Tahun Akademik 2018/2019.

\end{enumerate}

Penulis telah membuat laporan ini dengan sebaik-baiknya, diharapkan memberikan kritik dan saran dari semua pihak yang bersifat membangun, terimakasih.

\begin{raggedleft}

Bandung, 2019

Penulis

\end{raggedleft}

\end{acknowledgements}   % include an acknowledgements.tex file

\addcontentsline{toc}{chapter}{DAFTAR ISI}
\tableofcontents % generate and include a table of contents

\addcontentsline{toc}{chapter}{DAFTAR GAMBAR}
\listoffigures

\addcontentsline{toc}{chapter}{DAFTAR TABEL}
\listoftables           % generate and include a list of figures

\addcontentsline{toc}{chapter}{DAFTAR SIMBOL}
\begin{simbol}

\end{simbol}




\addcontentsline{toc}{chapter}{DAFTAR SINGKATAN}
\begin{singkatan}
\thispagestyle{plain}
\begin{table}[ht]
	\centering
	\begin{tabular}{|c|c|c|}
	\hline
		\textbf{No} & \textbf{Singkatan} & \textbf{Keterangan}   \\
	\hline
		1. & IRC       & \textit{Informatics Research Center}  \\
		2. & Dibuat Oleh & Maulyanda, Resti, dan Ikhsan \\
	\hline    
	\end{tabular}
\end{table}

\end{singkatan}

\end{romanpages}            % end roman page numbering

%now include the files of latex for each of the chapters etc
\chapter{PENDAHULUAN}

\section{Latar Belakang}

\par
Logistik dan Manajemen Rantai Pasok merupakan sebuah elemen penting untuk praktek operasi logistik berkelanjutan \cite{grant2017sustainable}, yang bertujuan untuk menganalisa hubungan antar proses logistik baik berupa \textit{continue}, \textit{discrate}, kombinasi sampai informasi yang terkait dari titik pasokan sampai ke titik permintaan untuk dapat memenuhi permintaan pelanggan \cite{mahendrawathi2015service}. Logistik dan rantai pasokan memainkan peran penting dalam keberhasilan keseluruhan manajemen ritel \cite{fernie2018logistics}. 
\par Kualitas dalam pelayanan logstik dan manajemen rantai pasok juga dapat mempengaruhi standar kesuksesan suatu perusahaan dalam memanjakan \textit{customer} \cite{pane2019implementasi}, maka dibutuhkan inovasi terbaru yang dapat meningkatkan daya saing dalam layanan kepada konsumen, termasuk dalam mengelola data gudang. Gudang umumnya mengacu pada puluhan rak tempat penyimpanan barang, dengan peralatan penanganan material yang sesuai untuk operasi gudang kargo masuk dan keluar yang dapat menggunakan banyak lahan, memakan banyak waktu, biaya dan mampu memaksa pekerja untuk dapat bekerja dengan intensitas tinggi \cite{zhou2016warehouse}. Penggunaan data \textit{warehouse} hampir dibutuhkan oleh semua organisasi maupun perusahan karena data \textit{warehouse} dapat memungkinkan integrasi berbagai macam jenis data dari berbagai macam aplikasi atau sisitem \cite{khan2016privacy}, pada penelitian ini pengunnan teknologi integrasi data \textit{warehouse} digunakan pada alat \textit{smart conveyor} dan sistem WMS yang dapat mengatur data pemasukan, pengeluaran dan data barang permintaan terbanyak pada sebuah gudang \cite{pane2018qualitative}. Penggunaan teknologi integrasi data antara alat dan sistem dapat memaksimalkan kinerja dari \textit{warehouse managent system}. 
\par Permasalahan yang terdapat pada saat pengujian \textit{smart conveyor} dan sistem wms saat ini yaitu tidak ada sinkronisasi data yang dibaca oleh RFID dengan data yang masuk kedalam sistem wms, yang berarti data yang dibaca dengan hasil data \textit{inbound} yang ada pada sistem berbeda. Adanya kesalahan dalam mebaca data dapat menyebabkan kesalahan dalam perhitungan data pada sistem dan mengelola data dengan menggunakan algoritma K-Means dan metode \textit{moving average} \cite{awangga2018k} untuk penerapan perhitungan permintaan terbanyak.
\par Penyesesuaian intergrasi data antara data yang dibaca oleh \textit{smart conveyor} dengan data yang masuk kedalam sistem WMS dapat membenahi permasalahan dalam perhitungan data dan penerapan perhitungan permintaan terbanyak. Sehingga tidak ada lagi terjadi kesalahan dalam mengelola data gudang pada sistem \textit{warehouse management system}. \break


\section{Problems}
the problem you want to solve

\section{Objective and Contribution}
\subsection{Objective}
What your research purpose

\subsection{Contribution}
Whats your contribution

\section{Scoop and Environtment}
Scoop and environment for the research 
\chapter{LANDASAN TEORI}

Your related works, and your purpose and contribution which must be different as below.

\section{Same Topics}
Cite every latest journal with same topic
\subsection{Topic 1}
cite for first topic

\subsection{Topic 2}
if you have two topics you can include here to


\section{Same Method}
write and cite latest journal with same method

\subsection{Method 1}
cite and paraphrase method 1

\subsection{Method 2}
cite and paraphrase method 2 if you have more method please add new subsection. 
\chapter{GAMBARAN OBYEK STUDI}

\subsection{Data Primer}

\subsection{Data Sekunder}
 
\chapter{METODOLOGI PENELITIAN} 
\chapter{PENGUJIAN DAN HASIL} 
\chapter{KESIMPULAN}

\section{Kesimpulan Masalah}

\section{Kesimpulan Metode}

\section{Kesimpulan Pengujian Sistem} 
%\chapter{DISKUSI}

%next line adds the Bibliography to the contents page
\addcontentsline{toc}{chapter}{Daftar Pustaka}
%uncomment next line to change bibliography name to references
%\renewcommand{\bibname}{References}
\bibliography{references}        %use a bibtex bibliography file refs.bib
\bibliographystyle{plain}  %use the plain bibliography style

%now enable appendix numbering format and include any appendices
\appendix
\chapter{\textit{CURRICULUM VITAE}}


\chapter{JURNAL}


\end{document}

